\documentclass{wissdoc}
% ----------------------------------------------------------------
% Abschlussarbeit - Hauptdokument
% ----------------------------------------------------------------
%%
%%
% wissdoc Optionen: draft, relaxed, pdf --> siehe wissdoc.cls
% ------------------------------------------------------------------
% Weitere packages: (Dokumentation dazu durch "latex <package>.dtx")
\usepackage[numbers,sort&compress]{natbib}

\usepackage{ngerman} % deutsche Spracheinstellung für Umlaute etc., keine Rechtschreibprüfung
\usepackage[utf8]{inputenc} % korrekte Darstellung von Umlauten
\usepackage[printonlyused]{acronym} % Abkürzungsverzeichnis nur mit Abkürzungen, die auch verwendet werden

% useful packages:
% \usepackage{amsmath} % Standardpaket für Bereich Mathematik
% \usepackage{rotating} % Rotation von z.B. Bildern, Tabellen etc.
% \usepackage{enumitem} % Kontrolle über Layout von itemize, enumerate, description
% \usepackage{pdfpages} % Seiten aus PDF/ ganzes PDF einbinden
% \usepackage{newunicodechar} % Bedeutung von Unicode Characters definieren, Nutzung über \newunicodechar{}{}


% \usepackage{varioref}
% \usepackage{verbatim}
% \usepackage{float}    %z.B. \floatstyle{ruled}\restylefloat{figure}
% \usepackage{subfigure}
% \usepackage{fancybox} % für schattierte,ovale Boxen etc.
% \usepackage{tabularx} % automatische Spaltenbreite
% \usepackage{supertab} % mehrseitige Tabellen
%% ---------------- end of usepackages -------------


%% Informationen für die PDF-Datei
\hypersetup{
 pdfauthor={N.N.},
 pdftitle={Not set}
 pdfsubject={Not set},
 pdfkeywords={Not set}
}

% Print URLs not in Typewriter Font
\def\UrlFont{\rm}

\newcommand{\blankpage}{% Leerseite ohne Seitennummer, nächste Seite rechts
 \clearpage{\pagestyle{empty}\cleardoublepage}
}

%% Einstellungen für das gesamte Dokument

% Trennhilfen
% Wichtig! 
% Im ngerman-paket sind zusätzlich folgende Trennhinweise enthalten:
% "- = zusätzliche Trennstelle
% "| = Vermeidung von Ligaturen und mögliche Trennung (bsp: Schaf"|fell)
% "~ = Bindestrich an dem keine Trennung erlaubt ist (bsp: bergauf und "~ab)
% "= = Bindestrich bei dem Worte vor und dahinter getrennt werden dürfen
% "" = Trennstelle ohne Erzeugung eines Trennstrichs (bsp: und/""oder)

% Trennhinweise fuer Woerter hier beschreiben
\hyphenation{
% Pro-to-koll-in-stan-zen
% Ma-na-ge-ment  Netz-werk-ele-men-ten
% Netz-werk Netz-werk-re-ser-vie-rung
% Netz-werk-adap-ter Fein-ju-stier-ung
% Da-ten-strom-spe-zi-fi-ka-tion Pa-ket-rumpf
% Kon-troll-in-stanz
}

% Index-Datei öffnen
\ifnotdraft{\makeindex}
%%%%%%%%%%%%%% includeonly %%%%%%%%%%%%%%%%%%%
% Es werden nur die Teile eingebunden, die hier 
% aufgefuehrt sind!
\includeonly{%
chapters/abstract_de, % Abstract Deutsch
abkuerzungen, %  
titelseite,%
erklaerung,% Ist in KA Pflicht für Abschlussarbeiten
chapters/einleitung,% Motivation, Zielsetzung, Gliederung
chapters/grundlagen,% Grundlagen 
chapters/analyse,   % Problembeschreibung (Detail) und Related Work
chapters/entwurf,   % Beschreibung der Problemlösung (Konzepte, allg. Architektur, ...)
chapters/implementierung,  % Beschreibung der Umsetzung/Implementierung
chapters/evaluation,      % Nachweis und Auswertung
chapters/zusammenfassung,  % Zusammenfassung der Ergebnisse und Ausblick
anhang
}
%%%%%%%%%%%%%%%%%%%%%%%%%%%%%%%%%%%%%%%%%%%%%%
\begin{document}

\frontmatter
\pagenumbering{roman}
\ifnotdraft{
 \include{titelseite}
 \blankpage % Leerseite auf Titelrückseite
 %
 % Die folgende Erklärung ist für Diplomarbeiten Pflicht
 % (siehe Prüfungsordnung), für Studienarbeiten nicht notwendig
 \include{erklaerung}
 \blankpage % Leerseite auf Erklärungsrückseite
}
%
%% *************** Hier geht's ab ****************
\include{chapters/abstract_de}
%% ++++++++++++++++++++++++++++++++++++++++++
%% Verzeichnisse
%% ++++++++++++++++++++++++++++++++++++++++++
\setcounter{tocdepth}{2}
\ifnotdraft{
{\parskip 0pt\tableofcontents} % toc bitte einzeilig
\blankpage
\listoffigures
\blankpage
\listoftables
\blankpage
\include{abkuerzungen}
}


%% ++++++++++++++++++++++++++++++++++++++++++
%% Hauptteil
%% ++++++++++++++++++++++++++++++++++++++++++
\graphicspath{{Bilder/}}

\mainmatter
\pagenumbering{arabic}
\include{chapters/einleitung}  % Einleitung
\include{chapters/grundlagen}  % Grundlagen
\include{chapters/analyse}     % Analyse
\include{chapters/entwurf}     % Entwurf
\include{chapters/implementierung}    % Implementierung
\include{chapters/evaluation}        % Evaluierung
\include{chapters/zusammenfassung}   % Zusammenfassung und Ausblick

%% ++++++++++++++++++++++++++++++++++++++++++
%% Literatur
%% ++++++++++++++++++++++++++++++++++++++++++
%  mit dem Befehl \nocite werden auch nicht 
%  zitierte Referenzen abgedruckt
\cleardoublepage
\phantomsection
\addcontentsline{toc}{chapter}{\bibname}
%%
\nocite{*} % nur angeben, wenn auch nicht im Text zitierte Quellen 
           % erscheinen sollen
\bibliographystyle{accessabbrv} % mit abgekürzten Vornamen der Autoren
%\bibliographystyle{gerplain} % abbrvnat unsrtnat
% \bibliographystyle{alpha}
\bibliography{thesis}

%% ++++++++++++++++++++++++++++++++++++++++++
%% Anhang
%% ++++++++++++++++++++++++++++++++++++++++++

\appendix
\include{anhang}
%\include{anhang_b}

\end{document}
%% end of file
